\documentclass[10pt,a4paper]{amsart}
\usepackage{amsmath,amssymb,amsthm,graphicx}
\usepackage{verbatim}
\newenvironment{reference}{\comment}{\endcomment}
%\newenvironment{reference}{}{}
\newenvironment{slogan}{\comment}{\endcomment}
\newenvironment{history}{\comment}{\endcomment}

% For commutative diagrams we use Xy-pic
\usepackage[all]{xy}

% We use 2cell for 2-commutative diagrams.
\xyoption{2cell}
% Theorem environments.
%
\theoremstyle{plain}
\newtheorem{theorem}[subsection]{Theorem}
\newtheorem{proposition}[subsection]{Proposition}
\newtheorem{lemma}[subsection]{Lemma}

\theoremstyle{definition}
\newtheorem{definition}[subsection]{Definition}
\newtheorem{example}[subsection]{Example}
\newtheorem{exercise}[subsection]{Exercise}
\newtheorem{situation}[subsection]{Situation}

\theoremstyle{remark}
\newtheorem{remark}[subsection]{Remark}
\newtheorem{remarks}[subsection]{Remarks}

\numberwithin{equation}{subsection}

% Macros
%
\def\lim{\mathop{\mathrm{lim}}\nolimits}
\def\colim{\mathop{\mathrm{colim}}\nolimits}
\def\Spec{\mathop{\mathrm{Spec}}}
\def\Hom{\mathop{\mathrm{Hom}}\nolimits}
\def\Ext{\mathop{\mathrm{Ext}}\nolimits}
\def\SheafHom{\mathop{\mathcal{H}\!\mathit{om}}\nolimits}
\def\SheafExt{\mathop{\mathcal{E}\!\mathit{xt}}\nolimits}
\def\Sch{\mathit{Sch}}
\def\Mor{\mathop{\mathrm{Mor}}\nolimits}
\def\Ob{\mathop{\mathrm{Ob}}\nolimits}
\def\Sh{\mathop{\mathit{Sh}}\nolimits}
\def\NL{\mathop{N\!L}\nolimits}
\def\CH{\mathop{\mathrm{CH}}\nolimits}
\def\proetale{{pro\text{-}\acute{e}tale}}
\def\etale{{\acute{e}tale}}
\def\QCoh{\mathit{QCoh}}
\def\Ker{\mathop{\mathrm{Ker}}}
\def\Im{\mathop{\mathrm{Im}}}
\def\Coker{\mathop{\mathrm{Coker}}}
\def\Coim{\mathop{\mathrm{Coim}}}

% Boxtimes
%
\DeclareMathSymbol{\boxtimes}{\mathbin}{AMSa}{"02}

%
% Macros for moduli stacks/spaces
%
\def\QCohstack{\mathcal{QC}\!\mathit{oh}}
\def\Cohstack{\mathcal{C}\!\mathit{oh}}
\def\Spacesstack{\mathcal{S}\!\mathit{paces}}
\def\Quotfunctor{\mathrm{Quot}}
\def\Hilbfunctor{\mathrm{Hilb}}
\def\Curvesstack{\mathcal{C}\!\mathit{urves}}
\def\Polarizedstack{\mathcal{P}\!\mathit{olarized}}
\def\Complexesstack{\mathcal{C}\!\mathit{omplexes}}
% \Pic is the operator that assigns to X its picard group, usage \Pic(X)
% \Picardstack_{X/B} denotes the Picard stack of X over B
% \Picardfunctor_{X/B} denotes the Picard functor of X over B
\def\Pic{\mathop{\mathrm{Pic}}\nolimits}
\def\Picardstack{\mathcal{P}\!\mathit{ic}}
\def\Picardfunctor{\mathrm{Pic}}
\def\Deformationcategory{\mathcal{D}\!\mathit{ef}}
\usepackage{cleveref}
\crefname{originallemma}{original lemma}{original lemmas}
\Crefname{originallemma}{Original lemma}{Original lemmas}
\makeatletter
\newlabel{lemma-base-change-invariance-of-relative-dimension}{{0DRN}{source}{}{sourcepointer-1}{}}
\newlabel{lemma-base-change-invariance-of-relative-dimension@cref}{{[originallemma][0][]0DRN}{[1][0][]source}{}{}{}}
\makeatother
\crefname{originalreference}{original reference}{original references}
\Crefname{originalreference}{Original reference}{Original references}
\makeatletter
\newlabel{spaces-properties-lemma-topology-points}{{03BX}{source}{}{sourcepointer-2}{}}
\newlabel{spaces-properties-lemma-topology-points@cref}{{[originalreference][0][]03BX}{[1][0][]source}{}{}{}}
\makeatother
\crefname{originalreference}{original reference}{original references}
\Crefname{originalreference}{Original reference}{Original references}
\makeatletter
\newlabel{spaces-morphisms-lemma-openness-bounded-dimension-fibres}{{04NT}{source}{}{sourcepointer-3}{}}
\newlabel{spaces-morphisms-lemma-openness-bounded-dimension-fibres@cref}{{[originalreference][0][]04NT}{[1][0][]source}{}{}{}}
\makeatother
\crefname{originalreference}{original reference}{original references}
\Crefname{originalreference}{Original reference}{Original references}
\makeatletter
\newlabel{morphisms-lemma-flat-finite-presentation-CM-fibres-relative-dimension}{{02NM}{source}{}{sourcepointer-4}{}}
\newlabel{morphisms-lemma-flat-finite-presentation-CM-fibres-relative-dimension@cref}{{[originalreference][0][]02NM}{[1][0][]source}{}{}{}}
\makeatother
\crefname{originalreference}{original reference}{original references}
\Crefname{originalreference}{Original reference}{Original references}
\makeatletter
\newlabel{morphisms-lemma-smooth-omega-finite-locally-free}{{02G1}{source}{}{sourcepointer-5}{}}
\newlabel{morphisms-lemma-smooth-omega-finite-locally-free@cref}{{[originalreference][0][]02G1}{[1][0][]source}{}{}{}}
\makeatother
\title{Relative dimension is upper semi-continuous, and locally constant for smooth morphisms}
\author{The Stacks Project Authors}
\date{stacks-project:stacks-geometry.tex:local-snapshot-20261012; GNU FDL 1.2 or later}
\begin{document}
\maketitle
\noindent\textit{Editorial note.} All delivered mathematical text is the author's own, copied verbatim from the pinned Stacks Project chapter file stacks-geometry.tex; no statement or proof step was rewritten, reordered or removed.\par\medskip
\noindent\textit{Editorial note.} Span recomputation by label pairing (LANE\_BRIEF 10.2): in this chapter the author's proof environment follows the statement environment as a standalone block; the spec covers the complete environment pair (statement span, proof span) as recorded in the pool candidate.\par\medskip
\noindent\textit{Editorial note.} Occupation scan: combined claim\_key + statement\_anchor + generic-line scan over all registration files (active and withdrawn) found 0 hits for stacks-geometry.tex (ge-occupied.json, 2026-10-14); 207 route=geometry / chapter-list field hits in focused-seeds are alias noise (route name and alphabetical chapter enumerations), excluded per the ladder alias-blind-spot caveat.\par\medskip
\ifdefined\hypertarget\hypertarget{sourcepointer-1}{}\fi
\noindent\textit{Source reference.} Original lemma 0DRN: sources/stacks-chow/tags.txt, source lines 17228--17228. Cross-chapter standard prerequisite cited by the author inside the counted span; it is used as a pointer and its own proof is not part of this item.\par\medskip
\ifdefined\hypertarget\hypertarget{sourcepointer-2}{}\fi
\noindent\textit{Source reference.} Original reference 03BX: sources/stacks-chow/tags.txt, source lines 4364--4364. Cross-chapter standard prerequisite cited by the author inside the counted span; it is used as a pointer and its own proof is not part of this item.\par\medskip
\ifdefined\hypertarget\hypertarget{sourcepointer-3}{}\fi
\noindent\textit{Source reference.} Original reference 04NT: sources/stacks-chow/tags.txt, source lines 6091--6091. Cross-chapter standard prerequisite cited by the author inside the counted span; it is used as a pointer and its own proof is not part of this item.\par\medskip
\ifdefined\hypertarget\hypertarget{sourcepointer-4}{}\fi
\noindent\textit{Source reference.} Original reference 02NM: sources/stacks-chow/tags.txt, source lines 3485--3485. Cross-chapter standard prerequisite cited by the author inside the counted span; it is used as a pointer and its own proof is not part of this item.\par\medskip
\ifdefined\hypertarget\hypertarget{sourcepointer-5}{}\fi
\noindent\textit{Source reference.} Original reference 02G1: sources/stacks-chow/tags.txt, source lines 3212--3212. Cross-chapter standard prerequisite cited by the author inside the counted span; it is used as a pointer and its own proof is not part of this item.\par\medskip
\begin{lemma}
\label{lemma-relative-dimension-is-semi-continuous}
Let $f: \mathcal{T} \to \mathcal{X}$ be a locally of finite type morphism of
algebraic stacks.
\begin{enumerate}
\item
The function $t \mapsto \dim_t(\mathcal{T}_{f(t)})$ is upper semi-continuous
on $|\mathcal{T}|$.
\item If $f$ is smooth, then
the function $t \mapsto \dim_t(\mathcal{T}_{f(t)})$ is locally constant
on $|\mathcal{T}|$.
\end{enumerate}
\end{lemma}

\begin{proof}
Suppose to begin with that $\mathcal{T}$ is a scheme $T$,
let $U \to \mathcal{X}$ be a smooth surjective morphism whose source
is a scheme, and let $T' = T \times_{\mathcal{X}} U$.
Let $f': T' \to U$ be the pull-back of $f$ over $U$,
and let $g: T' \to T$ be the projection.

\medskip\noindent
Lemma \ref{lemma-base-change-invariance-of-relative-dimension}
shows that $\dim_{t'}(T'_{f'(t')}) = \dim_{g(t')}(T_{f(g(t'))}),$
for $t' \in T'$, while,
since $g$ is smooth and surjective (being the base-change
of a smooth surjective morphism) the map induced by $g$ on underlying
topological spaces is continuous and open
(by
Properties of Spaces, Lemma \ref{spaces-properties-lemma-topology-points}), and
surjective. Thus it suffices to note that part (1) for the morphism $f'$
follows from
Morphisms of Spaces, Lemma
\ref{spaces-morphisms-lemma-openness-bounded-dimension-fibres}, and part (2)
from either of Morphisms, Lemma
\ref{morphisms-lemma-flat-finite-presentation-CM-fibres-relative-dimension}
or
Morphisms, Lemma \ref{morphisms-lemma-smooth-omega-finite-locally-free}
(each of which gives the result for schemes, from which
the analogous results for algebraic spaces can
be deduced exactly as in
Morphisms of Spaces, Lemma
\ref{spaces-morphisms-lemma-openness-bounded-dimension-fibres}.

\medskip\noindent
Now return to the general case,
and choose a smooth surjective morphism
$h:V \to \mathcal{T}$ whose source is a scheme.
If $v \in V$, then, essentially by definition,
we have
$$
\dim_{h(v)}(\mathcal{T}_{f(h(v))}) =
\dim_{v}(V_{f(h(v))}) - \dim_{v}(V_{h(v)}).
$$
Since $V$ is a scheme, we have proved that the first
of the terms on the right hand side of this equality
is upper semi-continuous (and even locally
constant if $f$ is smooth), while the second term is
in fact locally constant.
Thus their difference is upper semi-continuous
(and locally constant if $f$ is smooth),
and hence the function
$\dim_{h(v)}(\mathcal{T}_{f(h(v))})$
is upper semi-continuous on $|V|$ (and locally
constant if $f$ is smooth).
Since the morphism $|V| \to |\mathcal{T}|$ is open and surjective,
the lemma follows.
\end{proof}

\end{document}
